\section{Datengrundlage und synthetische Generierung}\label{sec:daten}

Die ursprüngliche Planung sah rund 2000 fotografierte und von Hand gelabelte
Kartenhaufen vor. Erste Versuche mit dem Online-Dienst Roboflow zeigten
schnell, dass weder das Fotografieren noch das Labeln dieser Menge realistisch
war. 2000 Bilder entsprechen bei 36 Klassen rund 56 Bildern je Klasse, also
3.7\,\% der von Ultralytics empfohlenen Untergrenze von 1500 Bildern je Klasse
\cite{ultralytics2025tips}. Der Entscheid fiel deshalb, die Karten einmal
zu fotografieren und alle Haufen daraus zu rendern. Dieses Kapitel
stellt die Quellscans und die unabhängige reale Prüfbasis mit ihren Labelregeln vor,
danach den Generator und seine Datensatzvarianten.

\subsection{Quellmaterial: 72 Kartenscans}\label{sec:scans}

Beide Blätter wurden Karte für Karte mit dem iPhone fotografiert, perspektivisch
begradigt und auf ein einheitliches Format gebracht: 1710$\times$2670 Pixel,
JPEG mit Qualität 75, ohne Farbprofil und Metadaten.

Durch wechselnde Lichtverhältnisse, insbesondere Sonneneinstrahlung, wiesen
einige Scans Farbunterschiede auf. Sie wurden soweit möglich an die übrigen
Karten angeglichen, damit das Modell den Farbstich nicht als Merkmal einer
einzelnen Karte lernt. Weitere Korrekturen beschreibt
Abschnitt~\ref{sec:label-qa}.

\subsection{Reales Validierungsset und Labeling-Werkzeug}\label{sec:labeling}

Das aktuelle Validierungsset umfasst 813 eigene Aufnahmen: 305 Karten des
französischen Blatts, 208 des deutschen Blatts und 300 Negative. Sie stammen
aus sieben Videos und einer Fotoserie des französischen Blatts, sechs Videos
des deutschen Blatts sowie Fotos ohne Karte. Die Aufnahmen zeigen
unterschiedliche Unterlagen, Lichtverhältnisse und Blickwinkel
(Abbildung~\ref{fig:realval}). Viele Frames stammen aus denselben Videos
und sind deshalb nicht unabhängig. Trainiert wird ausschliesslich auf
synthetischen Bildern.

\begin{figure}[H]
\centering
\includegraphics[width=0.85\textwidth]{real-val-beispiele.jpg}
\caption{Frames des realen Validierungssets mit ihren Labels: drei Aufnahmequellen, zwei flache Winkel und ein Negativ.}
\label{fig:realval}
\end{figure}

Das eigene Dataset Tool zeigt Bilder mit ihren Labels in allen Varianten.
Beim Labeln werden die vier Ecken der obersten Karte und ihre Klasse
gesetzt. Ein Frame ohne gültige Karte wird als \enquote{No card}
gespeichert. Aus einer Annotation entstehen alle vier Label-Formate.
Die Vorschau des entzerrten Ausschnitts zeigt, ob die Eckreihenfolge
stimmt (Abbildung~\ref{fig:dataset-tool-labeler}). Nicht quadratische
Fotos werden durch einen mittigen Ausschnitt auf 640$\times$640 gebracht.
Das entspricht dem Sucherquadrat der App. Das Werkzeug zerlegt Videos in
Frames und erlaubt Korrekturen an gespeicherten Labels. Die häufigen
Schritte lassen sich über die Tastatur bedienen.

\begin{figure}[H]
\centering
\includegraphics[width=0.9\textwidth]{dataset-tool-labeler.png}
\caption{Das eigene Dataset Tool beim Labeln eines realen Frames. Die vier Ecken markieren die oberste Karte. Rechts ist der entzerrte Ausschnitt zu sehen.}
\label{fig:dataset-tool-labeler}
\end{figure}

\subsection{Grenzfälle: was das Modell erkennen muss}\label{sec:grenzfaelle}

Die Grenze zwischen Karte und Negativ ist bei verwischten, teilweise
verdeckten oder sehr flach aufgenommenen Karten nicht immer eindeutig.
Für jedes reale Frame wurde entschieden, ob die oberste Karte lesbar ist.
Lesbare Karten wurden mit Klasse und Ecken gelabelt. Frames ohne lesbare
Karte wurden als \enquote{No card} gespeichert, unklare Fälle teilweise
ausgelassen. Das bleibt ein Problem für die Auswertung: Erkennt das Modell
eine Karte in einem als negativ gelabelten Frame, zählt dies als
Fehlalarm. Wird
eine schwierig lesbare, aber gelabelte Karte übersehen, zählt dies als
Fehler. Die Einteilung ist für das bestehende Set festgelegt. Ihre
Wirkung auf Fehlalarme muss mit weiteren Aufnahmen geprüft werden.
Eine eigene Markierung für schwierige Frames und eine getrennte
Auswertung sind noch nicht umgesetzt (Abschnitt~\ref{sec:ausblick}).

\subsection{Synthetische Szenen}\label{sec:szenen}

Eine Szene wird in drei Ebenen von hinten nach vorn aufgebaut. Zuunterst
liegt der \emph{Tischhintergrund}: ein prozedural erzeugtes, gleichmässig
wiederholtes Muster (Uni, Filz, Sprenkel, Karo, Tartan) in zufälliger Farbe.
Es liegt als Textur auf der Tischebene und wird mit derselben Kamera
projiziert wie die Karten, damit Fluchtpunkt und Perspektive übereinstimmen.
Darauf werden 0 bis 35 \emph{Unterkarten} aus dem Rest desselben Blatts
um die Mitte gestreut und gedreht. Zuoberst liegt die \emph{oberste Karte},
die allein gelabelt wird. Ihre Position variiert, damit ein Klassifikator
nicht \enquote{Mitte gleich Antwort} lernt.

Die Kamera ist eine Lochkamera, die aus einer zufälligen Position über dem
Tisch auf einen Punkt nahe der Mitte blickt. Für jede Karte werden die vier
Weltecken projiziert und der Kartenscan per Homographie exakt in das
projizierte Viereck verzerrt \cite{opencv2025calib3d}. Die Kameradistanz
variiert, die Neigung reicht von 0\textdegree{} bis 60\textdegree{} und
die Belichtung schwankt um $\pm$20\,\%. Jede Karte liegt 0.3\,mm höher als
die darunterliegende. Bei einem vollen Blatt wird der Stapel dadurch
rund einen Zentimeter hoch. Die abgerundeten Kartenecken bleiben
transparent und folgen der Perspektive.

In rund 65\,\% der Szenen beleuchtet eine Lampe den Tisch. Die übrigen
zeigen diffuses Raum- oder Tageslicht. Ein Teil der Lampenszenen erhält
eine zweite, aufhellende Lichtquelle. Das Lichtzentrum kann auch neben
oder hinter dem Haufen liegen. Die Karten werfen weiche Kontaktschatten,
und eine Vignette dunkelt die Bildecken ab. Abbildung~\ref{fig:synthetisch}
zeigt zwölf Beispiele.

\begin{figure}[H]
\centering
\includegraphics[width=0.85\textwidth]{synthetisch-beispiele.jpg}
\caption{Gerenderte Trainingsbilder beider Blätter mit der Box der obersten Karte (Variante~C). Der Winkel ist die Kameraneigung gegenüber der Senkrechten.}
\label{fig:synthetisch}
\end{figure}

\subsection{Kamerawinkel und Blatt aus dem Bildindex}\label{sec:sweep}

Bildindex $i$ und Bildanzahl $n$ legen die Kameraneigung und das
Kartenblatt fest. Die Neigung von Bild $i$ beträgt
\[
  \theta_i = 60^\circ \cdot \frac{p(i) + \tfrac{1}{2}}{n},
\]
wobei $p$ die Bildindizes permutiert. Die Bilder decken so den Bereich
von 0\textdegree{} bis 60\textdegree{} gleichmässig ab, ohne nach Winkel
sortiert zu sein. Bei grösseren Neigungen kann der Horizont ins Bild
geraten und die Tischprojektion fehlerhaft werden.

Gerade Positionen von $p(i)$ erhalten das französische Blatt, ungerade
das deutsche. Bei den verwendeten geraden Bildzahlen bekommt jedes Blatt
genau die Hälfte der Bilder und einen eigenen gleichmässigen Winkelsweep.
Ein Kartenhaufen enthält immer nur Karten desselben Blatts.

\subsection{Sichtbarkeitsregel}\label{sec:sichtbarkeit}

Eine Karte wird nur gelabelt, wenn sie vollständig im Bild liegt und
mindestens drei Pixel Abstand zum Rand hat. Sind höchstens 85\,\% ihres
projizierten Rechtecks sichtbar, gilt das Bild als Negativ. Szenen
dazwischen werden erneut gezogen. Bleibt der Fall nach 16 Versuchen
uneindeutig, entsteht ein Negativ. Die Wiederholung verringert
widersprüchliche Labels bei fast gleichen Bildern. Bei realen Fotos
werden klar
angeschnittene Karten als \enquote{No card} gespeichert. Berührt eine
fast vollständig sichtbare Karte nur den Rand, wird das Frame
ausgelassen.

\subsection{Negativbeispiele und Bewegungsunschärfe}\label{sec:negative}

Der Generator sieht für Negative einen Anteil von 30\,\% vor. Das liegt
bewusst über den 0 bis 10\,\% Hintergrundbildern, die Ultralytics
empfiehlt \cite{ultralytics2025tips}. Gerenderte
Negative zeigen einen leeren Tisch, eine stark verwischte geworfene Karte
oder eine verwischte Karte, die die oberste Karte verdeckt. Im Videostrom
der App kommen solche Frames häufig vor. Eine fälschlich gezählte Karte
wiegt dabei schwerer als eine übersehene. Die Detektoren erhalten für
Negative kein Label, der Klassifikator der Variante~A die Klasse
\code{none}.

Für Bewegungsunschärfe wird die Karte entlang ihrer Bewegungsrichtung
mehrfach gezeichnet. Bei 30\,\% der positiven Bilder beträgt die
Bewegungsstrecke 0.02 bis 0.09 Kartenhöhen. Diese Karten bleiben lesbar.
Stark verwischte Karten mit 0.25 bis 0.65 Kartenhöhen gehören zu den
Negativen. Der Bereich dazwischen wird nicht erzeugt. Ein Viertel der
positiven Bilder zeigt zudem eine verwischte Karte am Bildrand, während
die oberste Karte klar sichtbar bleibt. In 5\,\% der Szenen liegt statt eines
gestreuten Haufens ein bündig gestapeltes Deck.

Rund ein Viertel der Negative stammt aus 25 realen Fotos ohne Karte.
Sie zeigen unter anderem Papier, Kassenbons, Schachteln, Tische und
Räume. Zufällige Ausschnitte, Drehungen, Spiegelungen und wechselnde
Belichtung machen daraus viele Trainingsbilder. Diese Fotos ergänzen
die gerenderten Negative um Alltagsgegenstände, die sonst als Karte
fehlgedeutet werden könnten.

\subsection{Label-Varianten und Datensatz-Layout}\label{sec:layout}

Aus einem Renderlauf schreibt der Generator alle Label-Formate der drei
Varianten (Tabelle~\ref{tab:labelvarianten}). Die Bilder werden einmal
gerendert und von den Varianten gemeinsam genutzt.

\begin{table}[H]
\centering\small
\begin{tabularx}{\textwidth}{@{}llYl@{}}
\toprule
\textbf{Variante} & \textbf{Ordner} & \textbf{Label} & \textbf{Klassen} \\
\midrule
C & \code{detect/} & achsenparallele Box \code{klasse cx cy w h}, eine je Bild & 72 \\
B$_1$ & \code{locate/} & orientierte Box aus vier Eckpunkten & 1 \\
B$_2$ & \code{classify\_crop/} & entzerrter Kartenausschnitt in einem Klassenordner & 72 \\
A & \code{classify\_full/} & ganzes Bild in einem Klassenordner, Negative unter \code{none} & 73 \\
\bottomrule
\end{tabularx}
\caption{Die vier Label-Varianten aus einem Renderlauf.}
\label{tab:labelvarianten}
\end{table}

Für Variante~B$_2$ wird die oberste Karte anhand ihrer vier Ecken per
Homographie zu einem aufrechten Bild von 384 Pixeln Höhe entzerrt. Im
Betrieb geschieht dies anhand der Vorhersage der ersten Stufe. Innerhalb
eines Datensatzes bilden Bilder, Detektions- und Lokalisierungslabels
sowie die Klassenliste die Grundlage. Abgeleitete Klassenordner und
die \code{data.yaml} werden bei Bedarf mit \code{rebuild} rekonstruiert.

\subsection{Qualitätssicherung der Labels}\label{sec:label-qa}

Automatische Konsistenzprüfungen und die Sichtung aller entzerrten
Ausschnitte fanden drei falsche Klassenlabels unter 211 geprüften Bildern.
Sie wurden korrigiert. Bei einigen Kreuz- und Schaufelkarten begann die
Vier-Ecken-Annotation an der falschen Ecke. Das dreht nur den entzerrten
Ausschnitt für B$_2$ um 180\textdegree{}. Die Boxen von C und B$_1$ bleiben
gleich. Automatische Korrekturversuche erzeugten zu viele Fehlmarkierungen,
daher wurden die Ausschnitte visuell geprüft. Dabei fielen auch zwölf
verdrehte französische Kartenscans auf, die korrigiert wurden.
